
\documentclass{article}
\usepackage{colortbl}
\usepackage{makecell}
\usepackage{multirow}
\usepackage{supertabular}

\begin{document}

\newcounter{utterance}

\centering \large Interaction Transcript for game `cladder', experiment `full\_v1.5\_default', episode 4386 with refineonly{-}Qwen3.5{-}27B.
\vspace{24pt}

{ \footnotesize  \setcounter{utterance}{1}
\setlength{\tabcolsep}{0pt}
\begin{supertabular}{c@{$\;$}|p{.15\linewidth}@{}p{.15\linewidth}p{.15\linewidth}p{.15\linewidth}p{.15\linewidth}p{.15\linewidth}}
    \# & \multicolumn{2}{c}{Player} && \multicolumn{2}{c}{Game Master} \\
    \hline

    \theutterance \stepcounter{utterance}  
    & & & \multicolumn{4}{p{0.6\linewidth}}{
        \cellcolor[rgb]{0.9,0.9,0.9}{
            \makecell[{{p{\linewidth}}}]{
                \texttt{\tiny{[P1$\langle$GM]}}
                \texttt{You are an expert in causal inference. The following question is not a typical commonsense query, but rather a meticulously designed question created by a professor specializing in causal inference, intended to assess the students' mastery of the course content.} \\
\texttt{Imagine a self{-}contained, hypothetical world with only the following conditions, and without any unmentioned factors or causal relationships: Gender has a direct effect on smoking and peanut allergy. Smoking has a direct effect on tar deposit. Tar deposit has a direct effect on peanut allergy. Gender is unobserved. For nonsmokers and with no tar deposit, the probability of being allergic to peanuts is 35\%. For nonsmokers and with high tar deposit, the probability of being allergic to peanuts is 63\%. For smokers and with no tar deposit, the probability of being allergic to peanuts is 25\%. For smokers and with high tar deposit, the probability of being allergic to peanuts is 49\%. For nonsmokers, the probability of high tar deposit is 22\%. For smokers, the probability of high tar deposit is 73\%. For smokers, would it be more likely to see being allergic to peanuts if the person had been a nonsmoker?} \\
\texttt{Start your answer with 'Yes' or 'No', followed by additional reasoning or evidence to support your explanation.} \\
            }
        }
    }
     \\ \\

    \theutterance \stepcounter{utterance}  
    & \multicolumn{4}{p{0.6\linewidth}}{
        \cellcolor[rgb]{0.561,0.737,0.561}{
            \makecell[{{p{\linewidth}}}]{
                \texttt{\tiny{[P1$\rangle$GM]}}
                \texttt{No} \\
\\ 
\texttt{To determine whether a smoker would be more likely to be allergic to peanuts if they had been a nonsmoker, we must compare the observed probability of allergy for smokers, \$P(Y=1 | S=1)\$, with the counterfactual probability of allergy for smokers had they not smoked, \$P(Y=1 | \textbackslash{}text\{do\}(S=0))\$.} \\
\\ 
\texttt{**1. Causal Structure and Identification**} \\
\texttt{The causal graph is defined as:} \\
\texttt{*   \$G \textbackslash{}to S\$ (Gender affects Smoking)} \\
\texttt{*   \$G \textbackslash{}to Y\$ (Gender affects Peanut Allergy)} \\
\texttt{*   \$S \textbackslash{}to T\$ (Smoking affects Tar Deposit)} \\
\texttt{*   \$T \textbackslash{}to Y\$ (Tar Deposit affects Peanut Allergy)} \\
\texttt{*   \$S \textbackslash{}to Y\$ (Direct effect of Smoking on Allergy, implied by the data structure where \$P(Y|S,T)\$ varies by \$S\$)} \\
\\ 
\texttt{Since \$G\$ is unobserved, we cannot simply compare \$P(Y|S=1)\$ and \$P(Y|S=0)\$ to find the causal effect. However, the problem provides the conditional probabilities \$P(Y|S, T)\$ and \$P(T|S)\$. This allows us to calculate the counterfactual outcome for the smoker population under the intervention \$\textbackslash{}text\{do\}(S=0)\$ using the **Backdoor Adjustment** or simply by integrating over the mediator \$T\$ using the specific conditional distributions provided for the intervention.} \\
\\ 
\texttt{The counterfactual probability \$P(Y=1 | \textbackslash{}text\{do\}(S=0))\$ for the population of smokers is calculated by assuming the smoking status is set to 0, while the distribution of the mediator \$T\$ follows the mechanism \$P(T|S=0)\$, and the outcome \$Y\$ follows \$P(Y|S=0, T)\$.} \\
\\ 
\texttt{**2. Calculate Observed Probability for Smokers (\$P(Y=1 | S=1)\$)**} \\
\texttt{Using the Law of Total Probability conditioning on Tar Deposit (\$T\$):} \\
\texttt{*   \$P(T=1 | S=1) = 0.73\$} \\
\texttt{*   \$P(T=0 | S=1) = 1 {-} 0.73 = 0.27\$} \\
\texttt{*   \$P(Y=1 | S=1, T=1) = 0.49\$} \\
\texttt{*   \$P(Y=1 | S=1, T=0) = 0.25\$} \\
\\ 
\texttt{\$\$P(Y=1 | S=1) = (0.49 \textbackslash{}times 0.73) + (0.25 \textbackslash{}times 0.27)\$\$} \\
\texttt{\$\$P(Y=1 | S=1) = 0.3577 + 0.0675 = 0.4252\$\$} \\
            }
        }
    }
    & & \\ \\

    \theutterance \stepcounter{utterance}  
    & & & \multicolumn{2}{p{0.3\linewidth}}{
        \cellcolor[rgb]{0.9,0.9,0.9}{
            \makecell[{{p{\linewidth}}}]{
                \texttt{\tiny{[GM$|$GM]}}
                \texttt{No} \\
            }
        }
    }
    & & \\ \\

    \theutterance \stepcounter{utterance}  
    & & & \multicolumn{2}{p{0.3\linewidth}}{
        \cellcolor[rgb]{0.9,0.9,0.9}{
            \makecell[{{p{\linewidth}}}]{
                \texttt{\tiny{[GM$|$GM]}}
                \texttt{no} \\
            }
        }
    }
    & & \\ \\

    \theutterance \stepcounter{utterance}  
    & & & \multicolumn{2}{p{0.3\linewidth}}{
        \cellcolor[rgb]{0.9,0.9,0.9}{
            \makecell[{{p{\linewidth}}}]{
                \texttt{\tiny{[GM$|$GM]}}
                \texttt{game\_result = WIN} \\
            }
        }
    }
    & & \\ \\

\end{supertabular}
}

\end{document}
